\documentclass[11pt]{article}
\usepackage[T1]{fontenc}
\usepackage[utf8]{inputenc}
\usepackage{lmodern,amsmath,amssymb,amsthm,mathtools,booktabs}
\usepackage[margin=1in]{geometry}
\usepackage{microtype}
\usepackage[colorlinks=true,linkcolor=blue,urlcolor=blue,citecolor=blue]{hyperref}
\usepackage{enumitem}
\hypersetup{pdftitle={Full support signed permutations and their asymptotic coefficients},pdfsubject={OEIS A260952 and uniform signed permutation asymptotics},pdfauthor={}}
\setlist{itemsep=3pt,topsep=5pt}
\newtheorem{theorem}{Theorem}
\newtheorem{lemma}[theorem]{Lemma}
\newtheorem{proposition}[theorem]{Proposition}
\newtheorem{corollary}[theorem]{Corollary}
\newcommand{\fall}[2]{(#1)_{#2}}
\newcommand{\coef}[2]{[#1^{#2}]}
\newcommand{\Stir}[2]{\genfrac\{\}{0pt}{}{#1}{#2}}
\newcommand{\N}{\mathbb N}
\DeclareMathOperator{\negnum}{neg}
\newenvironment{writenote}[1][]{\par\smallskip\begingroup\small
 \noindent\textbf{[write]}\if\relax\detokenize{#1}\relax\else~\textbf{#1.}\fi\ }%
 {\par\endgroup\smallskip}
\title{Full support signed permutations\\and their asymptotic coefficients}
\author{}
\date{2 October 2026}
\begin{document}
\maketitle
\begin{abstract}
We prove the explicit large-order conjecture recorded in OEIS A260952 and derive every fixed correction order. The coefficients belong to the factorial asymptotics of full-support elements of Weyl groups of types $B$ and $D$. An exact positive Stirling transform identifies the dominant scale $\log 3$ and gives a proof with controlled remainders. For the original signed-permutation model, two exact endpoint sums admit separate expansions uniformly near the zero-sign-weight boundary. They yield an all-order sparse-sign crossover and a Poisson law conditioned to be positive. Growth-index inverses are provided with an explicit rounding qualification. The underlying generating functions, connected-permutation asymptotics, and general Stirling-transform machinery are prior work; the contribution is the specialized conjecture resolution, its refinements, and the uniform crossover.
\end{abstract}

\section{The conjecture and the normalization}\label{fss:sec:conjecture}
Let
\begin{equation}\label{fss:eq:defs}
 g(x)=\sum_{n\ge0}n!x^n,\qquad C(x)=1-\frac1{g(x)}=\sum_{n\ge1}c_nx^n,
 \qquad B_n(q)=[x^n]\frac{g(qx)}{g(x)}.
\end{equation}
These are formal power series. The integers $c_n$ count indecomposable permutations and form OEIS \href{https://oeis.org/A003319}{A003319}. Put $b_n=B_n(2)$, which is \href{https://oeis.org/A109253}{A109253}, the number of elements of the Weyl group $B_n$ whose support contains every standard simple reflection. For type $D$, the corresponding sequence \href{https://oeis.org/A112225}{A112225}, denoted by $d_n$, satisfies the classical identity
\begin{equation}\label{fss:eq:D}
 d_n=\frac{b_n-3c_n}{2},\qquad n\ge2.
\end{equation}
Bergeron, Hohlweg, and Zabrocki~\cite{BHZ} give the generating functions and their combinatorial interpretation. Martin and Kearney~\cite{MK} give integral representations.

For fixed $q>1$, define $p_0(q)=1$ and the asymptotic coefficients $p_m(q)$ by
\begin{equation}\label{fss:eq:pdef}
 \frac{B_n(q)}{q^n n!}\sim\sum_{m\ge0}\frac{p_m(q)}{n^m}.
\end{equation}
The sequence $a_m=\href{https://oeis.org/A260952}{\mathrm{A260952}}(m)$ uses powers of $1/(2n)$ instead. Consequently
\begin{equation}\label{fss:eq:scaling}
 a_m=2^m p_m(2).
\end{equation}
Kot\v e\v sovec's 2015 conjecture, as stated in the OEIS entry, is
\begin{equation}\label{fss:eq:conj}
 a_m\sim-\frac{2^{m+1}m!}{9(\log3)^{m+1}}.
\end{equation}
The factor $2^m$ in~\eqref{fss:eq:scaling} must be retained when comparing the two normalizations.

\begin{theorem}\label{fss:thm:main}
The conjecture~\eqref{fss:eq:conj} holds. More precisely, with $\tau=\log3$ and $\fall mr=m(m-1)\cdots(m-r+1)$,
\begin{align}\label{fss:eq:main}
 a_m={}&-\frac{2^{m+1}m!}{9\tau^{m+1}}
 \left\{1-\frac{2\tau}{m}-\frac{9\tau^2}{4\fall m2}
 -\frac{27\tau^3}{2\fall m3}-\frac{1755\tau^4}{16\fall m4}
 +O(m^{-5})\right\}.
\end{align}
Every fixed subsequent order is computable by a finite formula with the corresponding next-order remainder. Moreover, $a_m<0$ for every $m\ge1$.
\end{theorem}
The next two coefficients in this falling-factorial form are
\[
 -\frac{9099\tau^5}{8\fall m5},\qquad
 -\frac{907821\tau^6}{64\fall m6}.
\]
They are obtained symbolically, rather than fitted to numerical data. The expansion is not asserted to converge.

\begin{writenote}[Added 2 October 2026, batch~77]
This report is a single manuscript, printed in full: manuscript~02 of
batch~77 of ProveIt's incoming-reports intake, delivered as the archive
\texttt{Coxeter\_\allowbreak Full\_\allowbreak Support\_\allowbreak Asymptotics.zip} (wrapper directory
\texttt{Coxeter\_\allowbreak Full\_\allowbreak Support/}, main file
\texttt{Coxeter\_\allowbreak Full\_\allowbreak Support\_\allowbreak Asymptotics.tex}, now \texttt{article.tex})
in commit \texttt{096ee7b87}, and placed in commit \texttt{d0e6008d9} in
the research-report collection, among the OEIS-sequence asymptotics
reports. The manuscript names no ProveIt commit and describes no search of
the repository, so it has no pin; its author line is empty. The bounded
source search of Section~\ref{fss:sec:repro} is the manuscript's own
literature search; the intake's untruncated search of the repository (the
research-report collection, the transseries volumes and \texttt{Oeis/})
found no other treatment of A260952, A109253 or A112225, or of full-support
elements of Coxeter groups. Nothing in this report is formalized in Lean or
Rocq, and placement in the collection confers no formal status on any
statement here. The writing step prefixed every label with \texttt{fss:},
added section labels and these \textsf{[write]} notes, and three
references \cite{fss-pdc,fss-tai,fss-cti}, and set the bibliography
ragged-right; no statement, proof or number of
the manuscript was altered. The type-$B$/$D$ models, generating functions
and connectivity characterization remain credited to \cite{BHZ}, and the
shifted-Stirling and compound-Poisson method to \cite{MK}, exactly as
delivered.

\emph{A neighbouring report.} The same batch delivered (as manuscript~17)
an all-orders study of the distinct parabolic double cosets of the
symmetric group (OEIS A260700), now the collection report
\texttt{a260700-\allowbreak parabolic-\allowbreak double-\allowbreak cosets}
\cite{fss-pdc}, which proves Browning's higher-order conjecture. Both
reports count objects defined by the standard parabolic subgroups of a
Coxeter group --- here the elements of the Weyl groups $B_n$ and $D_n$
lying in no proper standard parabolic subgroup, there the distinct double
cosets $W_IwW_J$ in type $A$ --- and both reach a factorial scale through
a Stirling transform and invert it with a Lambert $W$ seed. Their theorems,
generating functions and scales differ ($\tau=\log3$ and
$m!/(\log3)^{m+1}$ here; $\rho=\log2$ and $n!\rho^{-2n}$ there), and neither
report uses the other.

\emph{Notation.} Each symbol is printed as delivered. The following letters
carry two meanings here, or a different meaning in \cite{fss-pdc}.
\begin{center}\footnotesize
\begin{tabular}{@{}p{0.14\linewidth}p{0.80\linewidth}@{}}
\toprule
symbol & meaning here (and what it is \emph{not})\\
\midrule
$B_n$, $B_n(q)$ & the Weyl group of type $B$, and the sign-weighted full-support count \eqref{fss:eq:defs}, $b_n=B_n(2)$. Not the correction polynomials $B_m(r)$ of \cite{fss-pdc}.\\
$p_m(q)$, $a_m$ & $p_m(q)$ are \emph{asymptotic coefficients} \eqref{fss:eq:pdef}, $a_m=2^mp_m(2)$ is A260952 \eqref{fss:eq:scaling}; $a=(1+q)/q$ in Section~\ref{fss:sec:late}. In \cite{fss-pdc}, $p_n$ counts double cosets.\\
$c_n$ & indecomposable permutations (A003319); not the constant $c$ of Browning's conjecture in \cite{fss-pdc}.\\
$\tau$ & $\tau=\log3$ in Theorem~\ref{fss:thm:main}; $\tau=\log(1+q)$ in Section~\ref{fss:sec:late} and in $\tau(v)=\log(1+q/v)$; the same at $q=2$.\\
$D$ & type $D$ with $d_n$ (A112225); $D(n)=c_n/n!$ \eqref{fss:eq:Dseries} and its formal version $D(x)$; $D_0=\log(\alpha n_0)$ in Section~\ref{fss:sec:inverse}.\\
$L$, $H$, $\mathcal H$ & the endpoint sums $L_n(q)$, $H_n(q)$ and their generators $L(q,x)$, $H(q,x)$; $\mathcal H(t)$ is the binary entropy.\\
$K$, $K_n$ & the carrier constant $K$ of \eqref{fss:eq:carrier}; $K_n$ the number of negative signs (a random variable). Neither is the $K$ of \cite{fss-pdc}.\\
$A_1$, $A_2$, $A_0$ & pole coefficients in Section~\ref{fss:sec:late}; $A_0=\tfrac12\log(2\pi n_0)+\log K$ in Section~\ref{fss:sec:inverse}.\\
$u$, $u_j$; $v$, $v_j$ & $u$ is the negative-sign weight and $u_j$ the relative corrections of \eqref{fss:eq:carrier}; $v$ is a marking variable in the proof of Theorem~\ref{fss:thm:late} and $v_j$ the inverse coefficients of \eqref{fss:eq:invformal}.\\
$w$, $W$ & a signed permutation; $w=(e^t-1)/q$ in the proof of Theorem~\ref{fss:thm:late}; $W$ is the positive Lambert function.\\
$\ell$, $x$ & $\ell$ indexes $\gamma_\ell$ and is $\log y$ in Section~\ref{fss:sec:inverse}; $x$ is the generating-function variable and a formal proxy for $1/n$.\\
$\kappa(q)$, $\lambda$ & the limiting fraction in the proof of Theorem~\ref{fss:thm:late}; $\lambda$ is the sparse-sign parameter (not the $\lambda$ of \cite{fss-cti}'s unbalanced core).\\
\bottomrule
\end{tabular}
\end{center}
\end{writenote}

\section{A natural weight for negative signs}\label{fss:sec:weight}
For a signed permutation $w$, let $\negnum(w)$ count its negative entries and define
\[
 Z_n(u)=\sum_{\substack{w\in B_n\\w\text{ has full support}}}u^{\negnum(w)}.
\]
The usual sign-weighted extension of the classical generating function is
\begin{equation}\label{fss:eq:weighted}
 Z_n(u)=B_n(1+u),\qquad n\ge1.
\end{equation}
Thus the parameter has a direct combinatorial meaning.

Here is an elementary verification. Decompose the underlying unsigned permutation into its indecomposable blocks. The signed permutation has full support precisely when its final unsigned block contains a negative entry. Indeed, a missing simple reflection corresponds to a proper unsigned block cut with every signed entry to its right positive; the sign-change reflection is missing when every entry is positive. This is also a direct consequence of the connectivity characterization in~\cite[Proposition 23]{BHZ}.

If the final unsigned block has size $k$, there are $(n-k)!c_k$ underlying permutations. Writing $q=1+u$, their allowed sign weight is $q^{n-k}(q^k-1)$. Therefore
\begin{align}\label{fss:eq:positive}
 B_n(q)&=\sum_{k=1}^n(n-k)!c_kq^{n-k}(q^k-1)\\
 &=q^nn!-\sum_{k=1}^nc_kq^{n-k}(n-k)! .\label{fss:eq:exact}
\end{align}
In the second line we used $\sum_{k=1}^nc_k(n-k)!=n!$. In particular,
\begin{equation}\label{fss:eq:zero}
 B_n(1)=0\quad(n\ge1),
\end{equation}
and $Z_n(u)$ has nonnegative integer coefficients.

\section{The factorial expansion for connected permutations}\label{fss:sec:factorial}
Define the formal coefficients
\begin{equation}\label{fss:eq:beta}
 \sum_{j\ge0}\beta_jx^j=\frac1{g(x)^2}=(1-C(x))^2.
\end{equation}
The first values are $1,-2,-1,-4,-19,-110,-745$.

\begin{lemma}\label{fss:lem:factorial}
For each fixed integer $R\ge0$,
\begin{equation}\label{fss:eq:factorial}
 \frac{c_n}{n!}=\sum_{j=0}^R\frac{\beta_j}{\fall nj}+O_R(n^{-R-1}).
\end{equation}
\end{lemma}
\begin{proof}
The exact recurrence
\[
 c_n=n!-\sum_{k=1}^{n-1}c_k(n-k)!
\]
first gives $0<c_n\le n!$ and $c_n/n!=1+O(n^{-1})$. Split its convolution into the two ends and the middle. For fixed $R$,
\begin{equation}\label{fss:eq:recipbin}
 \sum_{k=R+1}^{n-R-1}\binom nk^{-1}=O_R(n^{-R-1}).
\end{equation}
For example, up to a sufficiently small fixed fraction of $n$, consecutive reciprocal-binomial terms decrease geometrically; in the central interval they are exponentially small. Since $c_k\le k!$, this controls the middle part of the convolution.

An induction on $R$, using the already obtained lower-order expansions at the opposite endpoint, now determines the factorial-shift coefficients:
\[
 \beta_0=1,\qquad
 \beta_r=-c_r-\sum_{j=1}^rj!\beta_{r-j}\quad(r\ge1).
\]
Their generating function is $(1-C)/g=1/g^2$, giving~\eqref{fss:eq:beta} and the stated remainder.
\end{proof}
This factorial-composition phenomenon is classical; see Comtet~\cite{Comtet} and Martin--Kearney~\cite{MK}. The proof is included to make the error control self-contained. In ordinary inverse powers it gives
\begin{equation}\label{fss:eq:Dseries}
 D(n):=\frac{c_n}{n!}\sim1-\frac2n-\frac1{n^2}-\frac5{n^3}
 -\frac{32}{n^4}-\frac{253}{n^5}-\cdots.
\end{equation}

\section{Two exact endpoints and their uniform expansions}\label{fss:sec:endpoints}
Set $h=\lfloor n/2\rfloor$ and define the finite sums
\begin{equation}\label{fss:eq:endpoints}
 \begin{aligned}
 L_n(q)&=1-\sum_{k=1}^h\frac{c_k}{q^k\fall nk},\\
 H_n(q)&=\sum_{j=0}^{n-h-1}\frac{q^j j!D(n-j)}{\fall nj}.
 \end{aligned}
\end{equation}
Equation~\eqref{fss:eq:exact} becomes the exact identity
\begin{equation}\label{fss:eq:twosectors}
 \frac{B_n(q)}{q^n n!}=L_n(q)-q^{-n}H_n(q).
\end{equation}

\begin{theorem}\label{fss:thm:endpoint}
For every fixed $R$, each of $L_n(q)$ and $H_n(q)$ has an expansion through $n^{-R}$ with error $O(n^{-R-1})$, uniformly on compact subsets of
\[
 \frac14<|q|<4.
\]
Every fixed derivative in $q$ has the corresponding uniform expansion. In particular the result holds on $[1,Q]$ for each fixed $Q<4$.
\end{theorem}
\begin{proof}
The tail estimates reduce to sums of $Q^j/\binom nj$ for $j\le n/2$ and $Q<4$. Up to $\delta n$, choose $\delta>0$ so that the consecutive ratio is at most $1/2$. Beyond this point, write the exponential rate as
\[
 t\log Q-\mathcal H(t),\qquad
 \mathcal H(t)=-t\log t-(1-t)\log(1-t).
\]
It is a convex function of $t$. At a sufficiently small $\delta$ and at $1/2$ its values are negative, since $\frac12\log Q-\log2<0$. The intervening tail is exponentially small. Removing the first $R+1$ terms consequently leaves $O_R(n^{-R-1})$. For $L_n$ use $Q=\max|q|^{-1}<4$; for $H_n$ use $Q=\max|q|<4$. Fixed derivatives introduce only fixed powers of the summation index and do not affect these conclusions. Finally apply Lemma~\ref{fss:lem:factorial} in the finitely many surviving terms of $H_n$.
\end{proof}
With $x$ a formal proxy for $1/n$, define
\begin{align}
 D(x)&=\sum_{j\ge0}\frac{\beta_jx^j}{\prod_{r=0}^{j-1}(1-rx)},\label{fss:eq:Dformal}\\
 L(q,x)&=1-\sum_{k\ge1}\frac{c_kq^{-k}x^k}{\prod_{r=0}^{k-1}(1-rx)},\label{fss:eq:Lformal}\\
 H(q,x)&=\sum_{j\ge0}\frac{q^j j!x^j}{\prod_{r=0}^{j-1}(1-rx)}
 D\!\left(\frac{x}{1-jx}\right).\label{fss:eq:Hformal}
\end{align}
Each fixed coefficient involves only finitely many terms, so these are effective all-order generators. Their first coefficients are
\begin{align}
 L(q,x)={}&1-\frac{x}{q}-\frac{x^2}{q^2}
 -\frac{q+3}{q^3}x^3-\frac{q^2+9q+13}{q^4}x^4+\cdots,\label{fss:eq:Lcoeff}\\
 H(q,x)={}&1+(q-2)x+(2q^2-2q-1)x^2\notag\\
 &+(6q^3-2q^2-3q-5)x^3\notag\\
 &+(24q^4+6q^3-12q^2-9q-32)x^4+\cdots.\label{fss:eq:Hcoeff}
\end{align}

For an arbitrary fixed $q>1$, including $q\ge4$, the single main expansion~\eqref{fss:eq:pdef} remains valid. Indeed the omitted original terms with $k>n/2$ in~\eqref{fss:eq:exact}, normalized by $q^nn!$, total at most $nq^{-n/2}$. The first half is controlled by reciprocal-binomial tails. This observation does \emph{not} extend the separate secondary-endpoint bounds to $q\ge4$.

At $q=1$, equations~\eqref{fss:eq:zero} and~\eqref{fss:eq:twosectors} imply $L_n(1)=H_n(1)$ exactly, hence $L(1,x)=H(1,x)$ as formal series. This cancellation is essential in the sparse-sign regime.

\paragraph{Exponential-scale qualification.}
The exact finite sums in~\eqref{fss:eq:endpoints} define the two pieces in~\eqref{fss:eq:twosectors}. Replacing $L_n$ by any fixed algebraic truncation leaves a polynomially small error, which is larger than $q^{-n}$. Thus the formula is not a claim that a fixed truncation of the main series alone resolves an exponentially small remainder. No optimal-truncation or complete infinite-transseries assertion is made.

\section{The late coefficients as a positive Stirling transform}\label{fss:sec:stirling}
Using Stirling numbers of the second kind,
\[
 \frac1{\fall nk}=\sum_{m\ge k}\Stir{m-1}{k-1}n^{-m}.
\]
It follows from~\eqref{fss:eq:Lformal} that
\begin{equation}\label{fss:eq:stirling}
 p_m(q)=-\sum_{k=1}^m c_kq^{-k}\Stir{m-1}{k-1},\qquad m\ge1.
\end{equation}
This also defines the rational-function extension of $p_m$ to every $q\ne0$. It is this extension that is used below for $0<q\le1$; the original one-sector count asymptotic~\eqref{fss:eq:pdef} is asserted only for $q>1$. Since every summand in~\eqref{fss:eq:stirling} is nonnegative and at least one is positive, $p_m(q)<0$ for $q>0$.

If $C_e(z)=\sum_{k\ge1}c_kz^k/k!$, the exact shifted Borel transform in its disk of convergence is
\begin{equation}\label{fss:eq:borel}
 \sum_{m\ge1}\frac{p_m(q)t^{m-1}}{(m-1)!}
 =-\frac1q C_e'\!\left(\frac{e^t-1}{q}\right).
\end{equation}
The general shifted-Stirling and compound-Poisson construction is explicitly developed in~\cite{MK}; it is prior methodology.

\section{Every fixed late coefficient order}\label{fss:sec:late}
For $q>0$, put
\[
 \tau=\log(1+q),\quad a=\frac{1+q}{q},\quad
 A_2=-\frac{q}{(1+q)^2},\quad A_1=\frac{q+2}{(q+1)^2}.
\]
For $\ell\ge0$, define
\begin{equation}\label{fss:eq:gamma}
 \gamma_\ell(q)=[\delta^\ell]\left\{-\frac1q
 \sum_{j=2}^{\ell+2}\frac{(-1)^{j-1}\beta_j}{(j-2)!}
 \left[a(1-e^{-\delta})\right]^{j-2}\right\}.
\end{equation}

\begin{theorem}\label{fss:thm:late}
For every fixed $J\ge0$, as $m\to\infty$,
\begin{align}\label{fss:eq:latethm}
 p_m(q)={}& A_2m!\tau^{-m-1}+A_1(m-1)!\tau^{-m}\notag\\
 &+\sum_{\ell=0}^J\gamma_\ell(q)(-1)^{\ell+1}\ell!
 (m-\ell-2)!\tau^{\ell+1-m}\notag\\
 &+O\!\left((m-J-3)!\tau^{J+2-m}\right).
\end{align}
The remainder is uniform for $q$ in any fixed compact subinterval of $(0,\infty)$.
\end{theorem}
\begin{proof}
We give a finite positive-transform proof, avoiding any assumption about analytic continuation of the full function $C_e$.

First substitute the factorial-shift expansion of Lemma~\ref{fss:lem:factorial} into~\eqref{fss:eq:stirling}. For any selected $M$, the error is bounded by a constant times $k!k^{-M}$. Consider the positive comparison weights
\[
 k!q^{-k}\Stir{m-1}{k-1}.
\]
When additionally marked by $v^k$, their exponential generating function is
\begin{equation}\label{fss:eq:marked}
 \sum_{m\ge1}\frac{t^{m-1}}{(m-1)!}
 \sum_{k=1}^mk!(v/q)^k\Stir{m-1}{k-1}
 =\frac{v/q}{\bigl(1-(v/q)(e^t-1)\bigr)^2}.
\end{equation}
For positive $v$ its nearest singularity is the double pole
\[
 \tau(v)=\log(1+q/v);
\]
all other poles are $\tau(v)+2\pi i r$, $r\in\mathbb Z\setminus\{0\}$. Residue extraction is uniform for $v$ near $1$ and $q$ in a compact positive interval. Consequently the exponential scale is $\tau(v)^{-m}$. The usual exponential-moment bounds imply that $k$ lies in $[\eta m,(1-\eta)m]$ apart from an exponentially small fraction of the total weight, for some $\eta>0$. Its limiting fraction is
\[
 \kappa(q)=\frac{q}{(q+1)\log(1+q)}\in(0,1).
\]
For clarity, the lower-tail bound uses $v<1$ and the upper-tail bound uses $v>1$; choosing the thresholds strictly below and above $\kappa(q)$ makes their exponential rates negative. Compactness gives a uniform choice. It follows that the transformed factorial-approximation error is $O(m^{-M})$ relative to the comparison sum. The finitely many small-$k$ terms are negligible on the factorial scale.

Next evaluate each factorial-shift model exactly. For $j=0$ and $j=1$, the model generating functions are $(1-w)^{-2}$ and $(1-w)^{-1}$. For $j\ge2$,
\begin{equation}\label{fss:eq:model}
 \sum_{k\ge j}\frac{(k-j)!}{(k-1)!}w^{k-1}
 =\frac{(-1)^{j-1}}{(j-2)!}(1-w)^{j-2}\log(1-w)+P_j(w),
\end{equation}
where $P_j$ is a polynomial. This identity follows by repeated integration, starting with $-\log(1-w)$ at $j=2$.

Substitute $w=(e^t-1)/q$ and $\delta=\tau-t$. Then
\[
 1-w=a(1-e^{-\delta}),\qquad
 \log(1-w)=\log\delta+\text{a function analytic at }\delta=0.
\]
The double and simple pole coefficients are $A_2,A_1$, and the logarithmic coefficients are exactly~\eqref{fss:eq:gamma}. These finite model generating functions are explicit rational and logarithmic functions of $e^t$, whose nearest singularity is $t=\tau$. Coefficient transfer is applied only to these explicit functions. The analytic polynomial pieces become entire exponential polynomials and are negligible on the factorial scale.

Finally, for $m-1>\ell$,
\begin{equation}\label{fss:eq:logtransfer}
 (m-1)![t^{m-1}](\tau-t)^\ell\log(\tau-t)
 =(-1)^{\ell+1}\ell!(m-\ell-2)!\tau^{\ell+1-m}.
\end{equation}
The first omitted logarithmic term has the size of the remainder in~\eqref{fss:eq:latethm}. Taking $M\ge J+4$ in the positive-transform estimate makes its error no larger. This proves the theorem.
\end{proof}

The first explicit relative corrections are
\begin{align}\label{fss:eq:generallate}
 \frac{p_m(q)}{A_2m!\tau^{-m-1}}={}&1-\frac{(q+2)\tau}{qm}
 -\frac{(q+1)^2\tau^2}{q^2\fall m2}
 -\frac{4(q+1)^3\tau^3}{q^3\fall m3}\notag\\
 &-\frac{(q+1)^3(23q+19)\tau^4}{q^4\fall m4}
 +O_q(m^{-5}).
\end{align}
Putting $q=2$ and multiplying by $2^m$ proves Theorem~\ref{fss:thm:main}, including its normalization.

\section{Sparse signs and a conditional Poisson law}\label{fss:sec:sparse}
Fix $\lambda>0$ and give each negative sign weight $u=\lambda/n$. Write $x=1/n$ formally.

\begin{theorem}\label{fss:thm:sparse}
For every fixed $R$, locally uniformly for complex $\lambda$,
\begin{equation}\label{fss:eq:sparse}
 \frac{B_n(1+\lambda/n)}{n!}
 =\sum_{r=0}^R\frac{F_r(\lambda)}{n^r}+O_R(n^{-R-1}),
\end{equation}
where the all-order coefficient generator is
\begin{align}\label{fss:eq:Fgen}
 \sum_{r\ge0}F_r(\lambda)x^r
 ={}&e^\lambda\exp\!\left(\sum_{s\ge2}
 \frac{(-1)^{s+1}\lambda^s}{s}x^{s-1}\right)
 L(1+\lambda x,x)\notag\\
 &-H(1+\lambda x,x).
\end{align}
On a compact disk containing zero, the remainder in~\eqref{fss:eq:sparse} is in fact $O_R(|\lambda|n^{-R-1})$.
\end{theorem}
\begin{proof}
Apply the two separate endpoint expansions in Theorem~\ref{fss:thm:endpoint} in a fixed complex neighborhood of $q=1$, and expand $n\log(1+\lambda/n)$. This proves~\eqref{fss:eq:sparse} and~\eqref{fss:eq:Fgen}, uniformly on compact complex sets. Every $F_r(0)$ is zero because $L(1,x)=H(1,x)$. The holomorphic remainder also vanishes at zero. Cauchy bounds on a slightly larger compact disk therefore improve its bound to $O_R(|\lambda|n^{-R-1})$.
\end{proof}
The first coefficients are
\begin{align}
 F_0(\lambda)&=e^\lambda-1,\label{fss:eq:F0}\\
 F_1(\lambda)&=-\left(1+\frac{\lambda^2}{2}\right)e^\lambda+1,\label{fss:eq:F1}\\
 F_2(\lambda)&=e^\lambda\left(\frac{\lambda^4}{8}+\frac{\lambda^3}{3}
 +\frac{\lambda^2}{2}+\lambda-1\right)-\lambda+1.\label{fss:eq:F2}
\end{align}
After division by $e^\lambda-1$, the improved remainder gives a relative $O_R(n^{-R-1})$ bound uniformly for $0<\lambda\le\Lambda$.

Under the probability on full-support signed permutations proportional to $(\lambda/n)^{\negnum(w)}$, let $K_n$ count negative signs. Its probability generating function is
\begin{equation}\label{fss:eq:pgf}
 \mathbb E[z^{K_n}]=\frac{B_n(1+\lambda z/n)}{B_n(1+\lambda/n)}
 \longrightarrow\frac{e^{\lambda z}-1}{e^\lambda-1}.
\end{equation}
Thus $K_n$ converges in distribution to a Poisson random variable of mean $\lambda$ conditioned to be positive. In particular,
\[
 \mathbb E[K_n]\longrightarrow\frac{\lambda e^\lambda}{e^\lambda-1}.
\]
Dividing the two series from Theorem~\ref{fss:thm:sparse} gives every fixed correction order in~\eqref{fss:eq:pgf}, locally uniformly in $z$.

There is also an exact one-negative-sign check:
\begin{equation}\label{fss:eq:oneminus}
 B_n'(1)=c_{n+1}.
\end{equation}
Indeed the identity $x^2g'(x)+(x-1)g(x)=-1$ implies
\[
 \frac{xg'(x)}{g(x)}=\frac{C(x)}x-1.
\]
The first one-negative counts are $1,3,13,71,461$. Accordingly the derivatives $F_r'(0)$ reproduce the expansion of $c_{n+1}/(nn!)$.

\section{Growth index inversion}\label{fss:sec:inverse}
The expansions above give three classes of positive sequences of the form
\begin{equation}\label{fss:eq:carrier}
 X_n=K\alpha^n\Gamma(n+1)
 \left(1+\sum_{j=1}^R\frac{u_j}{n^j}+O(n^{-R-1})\right).
\end{equation}
The parameters are as follows:
\begin{itemize}
\item For $|a_n|$, $\alpha=2/\log3$ and $K=2/(9\log3)$, with $u_j$ generated by~\eqref{fss:eq:main}.
\item For $B_n(q)$ at fixed $q>1$, $\alpha=q$, $K=1$, and $u_j=p_j(q)$. For type $D$, $\alpha=2$, $K=1/2$, and the same algebraic coefficients $p_j(2)$ apply, with the exact correction~\eqref{fss:eq:D} on the smaller factorial scale.
\item For $B_n(1+\lambda/n)$ at fixed $\lambda>0$, $\alpha=1$, $K=e^\lambda-1$, and $u_j=F_j(\lambda)/(e^\lambda-1)$.
\end{itemize}
The sequences are eventually increasing, since their consecutive ratios are asymptotic to $\alpha(n+1)$.

For $\ell=\log y$, let $W$ be the positive Lambert function and put
\begin{equation}\label{fss:eq:inversebase}
 n_0=\frac{\ell}{W(\alpha\ell/e)},\qquad
 D_0=\log(\alpha n_0),\qquad
 A_0=\frac12\log(2\pi n_0)+\log K.
\end{equation}
The first two inverse correction steps are
\begin{align}\label{fss:eq:inverse}
 \nu(\ell)={}&n_0-\frac{A_0}{D_0}
 -\frac{A_0^2/(2D_0^2)-A_0/(2D_0)+v_1}{n_0D_0}
 +O\!\left(\frac{(\log n_0)^2}{n_0^2}\right),\\[-2pt]
 &\hspace{38mm}v_1=\frac1{12}+u_1.\notag
\end{align}
This follows by two Taylor--Newton steps applied to the logarithm of~\eqref{fss:eq:carrier}. Every subsequent fixed order is obtained by formal Newton iteration on
\begin{equation}\label{fss:eq:invformal}
 \ell=x\bigl(\log(\alpha x)-1\bigr)
 +\frac12\log(2\pi x)+\log K+\sum_{j\ge1}v_jx^{-j},
\end{equation}
where $v_j$ combines the Stirling log-Gamma coefficients with
$\log(1+\sum_{j\ge1}u_jx^{-j})$.

\paragraph{Integer thresholds.}
If a sufficiently deep smooth truncation defines the real solution $\nu_R$ with logarithmic error $O(n^{-R-1})$, the true threshold
\[
 N(y)=\min\{n\ge n_*:X_n\ge y\}
\]
obeys the enclosure
\begin{equation}\label{fss:eq:ceil}
 N(y)=\left\lceil\nu_R(\log y)
 +O\!\left(\frac{\nu_R^{-R-1}}{\log\nu_R}\right)\right\rceil.
\end{equation}
Here $n_*$ lies in the eventually increasing range. This notation describes a rounding ambiguity of the stated width. It does not assert an exact ceiling decision when $y$ is too close to a sequence value, nor assume a canonical analytic interpolation of the discrete sequence.

\begin{writenote}[Added 2 October 2026, batch~77: an instance of repository results]
The inversion of this section is an instance of the factorial-core
inversion of ProveIt's transseries volumes \cite{fss-tai,fss-cti}, and no
novelty is claimed for the method. The phase of \eqref{fss:eq:carrier} is
$x(\log x+\log\alpha-1)+O(\log x)$, so the seed \eqref{fss:eq:inversebase}
is the exact solution of $X(\kappa\log X+d)=L$ with $\kappa=1$,
$d=\log\alpha-1$ and $L=\ell$ in \texttt{p0:prop:factorial-core} of
\cite{fss-tai}: there $v=W_0\bigl((L/\kappa)e^{d/\kappa}\bigr)=W(\alpha\ell/e)$
and $X=e^{v-d/\kappa}=(e/\alpha)e^v=\ell/v=n_0$, and the core slope
$\kappa\log X+d+\kappa=\log(\alpha n_0)$ is the $D_0$ above. It is equally
the unbalanced core \texttt{t2:eq:unbalanced-core} of \cite{fss-cti}
($\kappa=1$, $\lambda=\log\alpha-1$, with $L=\log y$ in place of
$\log(y/C)$). In the variable $s=\log X+d$ that equation reads
$s+\log s=\log\ell+d$, which is \texttt{p0:thm:lambert-core} with
$a=b=1$; since $b>0$ its branch rule selects $W_0$, the positive branch
used here. The corrections \eqref{fss:eq:inverse} and the formal Newton
iteration on \eqref{fss:eq:invformal} are the reversion about an
admissible core, \texttt{p0:thm:core-reversion} of \cite{fss-tai} (in
\cite{fss-cti}, the scaled local reversion \texttt{t2:prop:scaled}); the
linear--logarithmic specialization \texttt{p0:thm:lambert-centered} and the
balanced inverse \texttt{t2:thm:balanced-inverse} do not apply directly,
because the phase is $x\log x$ rather than $ax+b\log x$. The enclosure
\eqref{fss:eq:ceil} is the separation condition, part~(2) of
\texttt{p0:thm:staircase}: the threshold is determined when the two
ceilings at the ends of the stated width agree, and not otherwise. Part~(1)
of that theorem (exact rounding) needs an admissible interpolation of the
sequence itself, which the smooth truncation $\nu_R$ is not; that is the
qualification stated above. The generic staircase arithmetic is formalized
in Lean as \texttt{Fabius.staircase\_ceil} and
\texttt{Fabius.staircase\_separation} in
\nolinkurl{Analysis/FabiusFunction/Lean/FabiusFunction/StaircaseInversion.lean};
those statements concern an arbitrary monotone function and give no formal
status to anything in this section. Specific to this report are the three
carriers of \eqref{fss:eq:carrier} with their coefficients $u_j$, which
come from Theorem~\ref{fss:thm:main}, from \eqref{fss:eq:pdef} with
\eqref{fss:eq:stirling}, and from Theorem~\ref{fss:thm:sparse}.
\end{writenote}

\section{Reproducibility and attribution}\label{fss:sec:repro}
The accompanying offline Python scripts use exact integer recurrences and symbolic series. The verification includes the first displayed OEIS terms, weighted signed-permutation enumeration through $n=6$, and large-order, endpoint, and sparse-sign numerical checks through $600$. The coefficient generator produces five sparse-sign correction orders and six displayed late factorial-basis corrections. The numerical tests are consistency checks; the proofs and uniform remainders are given above.

The following boundaries are important. The $B/D$ models and generating functions are established in~\cite{BHZ}; their leading factorial asymptotics and several corrections are already recorded in~\cite{OEISB,OEISD}. The connected-permutation expansion is classical, and~\cite{MK} contains the general shifted-Stirling/compound-Poisson construction and the relevant integral representations. These are credited foundations. The specific late-growth statement in~\cite{OEISLate}, its explicit all-fixed-order refinement, and the signed-weight crossover are the targets proved here. A bounded source search found no prior statement of this full specialized package; exhaustive historical priority is not asserted.

\begin{writenote}[Added 2 October 2026, batch~77]
In the collection the two scripts and the delivered \texttt{build.sh} and
\texttt{replay.sh} are in \texttt{code/}, and the three recorded outputs and
\texttt{requirements.txt} in \texttt{data/}, all byte-identical to the
delivery; the delivered PDF and the checksum ledger \texttt{SHA256SUMS}
(verified at intake and retired) are not shipped and survive in the
archive. \texttt{replay.sh} changes to its own directory and then calls
\texttt{code/derive.py}, so in the collection layout (where it lives in
\texttt{code/}) it does not find the scripts; \texttt{verify.py} writes
\texttt{verification.json} into the \texttt{data/} directory next to its
own parent, that is, over the shipped file; \texttt{build.sh} compiles the
delivered file name \texttt{Coxeter\_\allowbreak Full\_\allowbreak Support\_\allowbreak Asymptotics.tex},
which is shipped only as this \texttt{article.tex}. The report's
\texttt{README.md} gives a rerun on a scratch copy that leaves the shipped
files untouched.
\end{writenote}

\begingroup\raggedright
\begin{thebibliography}{9}
\bibitem{BHZ} N. Bergeron, C. Hohlweg, and M. Zabrocki, \emph{Posets related to the connectivity set of Coxeter groups}, J. Algebra \textbf{303} (2006), 831--846. \href{https://arxiv.org/abs/math/0509271v3}{arXiv:math/0509271v3}. In particular Propositions 23--26.
\bibitem{MK} R. J. Martin and M. J. Kearney, \emph{Integral representation of certain combinatorial recurrences}, Combinatorica \textbf{35} (2015), 309--315. \href{https://doi.org/10.1007/s00493-014-3183-3}{doi:10.1007/s00493-014-3183-3}. The author manuscript gives the compound-Poisson construction and the type $B/D$ integral representations.
\bibitem{Comtet} L. Comtet, \emph{Sur les coefficients de l'inverse de la s\'erie formelle $\sum n!t^n$}, C. R. Acad. Sci. Paris A \textbf{275} (1972), 569--572. Classical background; Lemma~\ref{fss:lem:factorial} is proved independently here.
\bibitem{OEISLate} V. Kot\v e\v sovec, \href{https://oeis.org/A260952}{OEIS A260952}, submitted 5 August 2015. The large-order conjecture and its $1/(2n)$ normalization. Accessed 2 October 2026.
\bibitem{OEISB} M. Zabrocki and contributors, \href{https://oeis.org/A109253}{OEIS A109253}. Type $B$ full-support counts and the posted asymptotic coefficients. Accessed 2 October 2026.
\bibitem{OEISD} M. Zabrocki and contributors, \href{https://oeis.org/A112225}{OEIS A112225}. Type $D$ full-support counts, with offset $2$. Accessed 2 October 2026.
\bibitem{fss-pdc} [write] ProveIt research-report collection, the A260700
report \emph{All fixed order asymptotics for distinct parabolic double
cosets} (batch~77, manuscript~17),
\nolinkurl{SetTheory/Cardinals/docs/reports/generating-functions-and-asymptotics/oeis-sequence-asymptotics/a260700-parabolic-double-cosets}.
\bibitem{fss-tai} [write] ProveIt, \emph{Transseries and inversion}, the
transseries-and-inversion volume (\texttt{p0:thm:lambert-core},
\texttt{p0:prop:factorial-core}, \texttt{p0:thm:lambert-centered},
\texttt{p0:thm:core-reversion}, \texttt{p0:thm:staircase}),
\nolinkurl{Analysis/Transseries/docs/series-and-transseries/Transseries_And_Inversion/transseries_and_inversion.tex}.
\bibitem{fss-cti} [write] ProveIt, \emph{Combinatorial transseries
inverses} (\texttt{t2:eq:unbalanced-core}, \texttt{t2:prop:scaled},
\texttt{t2:thm:balanced-inverse}),
\nolinkurl{Analysis/Transseries/docs/series-and-transseries/Combinatorial_Transseries_Inverses/Combinatorial_Transseries_Inverses.tex}.
\end{thebibliography}
\endgroup
\end{document}
